\begin{proof}
We first note that $\npsplit$ runs $\Split$ only a constant number of times since $\pnum$ is constant. Hence its running time is asymptotically the same as that of $\Split$, namely $O(\lambda^{k(1-\delta)})$.

Let $L=\lceil\log_2\pnum\rceil$ and $\bar{\pnum}=2^L$. Consider first the case $\instance\getsr\zkcdistzero{R}$. After $r$ levels of the binary splitting recursion, each output list contains $2^r$ instances. By repeatedly applying \cref{thm:split} and taking a union bound over the $2^r$ split calls used to pass from level $r$ to level $r+1$, for every output list after $L$ levels we have
\[
\tvd\left(
(\hat{\instance}_1,\ldots,\hat{\instance}_{\bar{\pnum}}),
\zkcdistzero{\sqrt[\bar{\pnum}]{R}}^{\bar{\pnum}}
\right)
\leq
\sum_{r=0}^{L-1}
2^r\binom{k}{2}4^{\binom{k}{2}}3\lambda^k/\sqrt[2^{r+1}]{R}.
\]
Since $2^{r+1}\leq \bar{\pnum}$ for every $r\leq L-1$, we have
$1/\sqrt[2^{r+1}]{R}\leq 1/\sqrt[\bar{\pnum}]{R}$. Therefore,
\[
\tvd\left(
(\hat{\instance}_1,\ldots,\hat{\instance}_{\bar{\pnum}}),
\zkcdistzero{\sqrt[\bar{\pnum}]{R}}^{\bar{\pnum}}
\right)
\leq
(2\bar{\pnum}-1)\binom{k}{2}4^{\binom{k}{2}}3\lambda^k/\sqrt[\bar{\pnum}]{R}.
\]
The actual $\pnum$-$\Split$ output discards the final $\bar{\pnum}-\pnum$ coordinates. By the data-processing property of total variation distance, this truncation cannot increase the distance. This gives the claimed zero-solution bound.

The one-solution case is identical except that at each level we follow one good child guaranteed by \cref{thm:split}. Thus there exists one output list whose full $\bar{\pnum}$-tuple is within the same distance of
$\zkcdistone{\sqrt[\bar{\pnum}]{R}}^{\bar{\pnum}}$. Truncating this list to its first $\pnum$ coordinates again cannot increase the distance, which proves the claimed one-solution bound.
\end{proof}
